
\documentclass{article}
\usepackage{colortbl}
\usepackage{makecell}
\usepackage{multirow}
\usepackage{supertabular}

\begin{document}

\newcounter{utterance}

\centering \large Interaction Transcript for game `cladder', experiment `full\_v1.5\_default', episode 2504 with Qwen3.5{-}9B.
\vspace{24pt}

{ \footnotesize  \setcounter{utterance}{1}
\setlength{\tabcolsep}{0pt}
\begin{supertabular}{c@{$\;$}|p{.15\linewidth}@{}p{.15\linewidth}p{.15\linewidth}p{.15\linewidth}p{.15\linewidth}p{.15\linewidth}}
    \# & \multicolumn{2}{c}{Player} && \multicolumn{2}{c}{Game Master} \\
    \hline

    \theutterance \stepcounter{utterance}  
    & & & \multicolumn{4}{p{0.6\linewidth}}{
        \cellcolor[rgb]{0.9,0.9,0.9}{
            \makecell[{{p{\linewidth}}}]{
                \texttt{\tiny{[P1$\langle$GM]}}
                \texttt{You are an expert in causal inference. The following question is not a typical commonsense query, but rather a meticulously designed question created by a professor specializing in causal inference, intended to assess the students' mastery of the course content.} \\
\texttt{Imagine a self{-}contained, hypothetical world with only the following conditions, and without any unmentioned factors or causal relationships: Kidney stone size has a direct effect on treatment and recovery. Treatment has a direct effect on recovery. For patients who have small kidney stones and not receiving treatment, the probability of recovery is 88\%. For patients who have small kidney stones and receiving treatment, the probability of recovery is 72\%. For patients who have large kidney stones and not receiving treatment, the probability of recovery is 16\%. For patients who have large kidney stones and receiving treatment, the probability of recovery is 4\%. The overall probability of large kidney stone is 41\%. For patients receiving treatment, would it be less likely to see recovery if the patient had received no treatment?} \\
\texttt{Start your answer with 'Yes' or 'No', followed by additional reasoning or evidence to support your explanation.} \\
            }
        }
    }
     \\ \\

    \theutterance \stepcounter{utterance}  
    & \multicolumn{4}{p{0.6\linewidth}}{
        \cellcolor[rgb]{0.561,0.737,0.561}{
            \makecell[{{p{\linewidth}}}]{
                \texttt{\tiny{[P1$\rangle$GM]}}
                \texttt{No} \\
\\ 
\texttt{**Reasoning:**} \\
\\ 
\texttt{To answer this question, we must determine if the **Counterfactual Probability** of recovery given treatment (\$P(Y=1 | T=1, do(T=0))\$) is lower than the **Observed Probability** of recovery given treatment (\$P(Y=1 | T=1)\$). In causal inference terms, we are checking if the Average Treatment Effect (ATE) for the treated group is positive.} \\
\\ 
\texttt{Let \$S\$ denote stone size (Small or Large) and \$T\$ denote treatment (1 = Treatment, 0 = No Treatment). Let \$Y\$ denote recovery.} \\
\texttt{The observed probabilities are given as:} \\
\texttt{*   \$P(Y=1 | S=\textbackslash{}text\{Small\}, T=1) = 0.72\$} \\
\texttt{*   \$P(Y=1 | S=\textbackslash{}text\{Large\}, T=1) = 0.04\$} \\
\\ 
\texttt{We are given the marginal probability of having a large stone: \$P(S=\textbackslash{}text\{Large\}) = 0.41\$.} \\
\texttt{Therefore, the probability of having a small stone is: \$P(S=\textbackslash{}text\{Small\}) = 1 {-} 0.41 = 0.59\$.} \\
\\ 
\texttt{**Step 1: Calculate the observed probability of recovery for patients receiving treatment.**} \\
\texttt{\$\$P(Y=1 | T=1) = P(Y=1 | S=\textbackslash{}text\{Small\}, T=1)P(S=\textbackslash{}text\{Small\}) + P(Y=1 | S=\textbackslash{}text\{Large\}, T=1)P(S=\textbackslash{}text\{Large\})\$\$} \\
\texttt{\$\$P(Y=1 | T=1) = (0.72 \textbackslash{}times 0.59) + (0.04 \textbackslash{}times 0.41)\$\$} \\
\texttt{\$\$P(Y=1 | T=1) = 0.4248 + 0.0164 = 0.4412\$\$} \\
\\ 
\texttt{**Step 2: Calculate the counterfactual probability of recovery if treated patients had received no treatment.**} \\
\texttt{Since the problem states that "Kidney stone size has a direct effect on treatment and recovery" and "Treatment has a direct effect on recovery," and no other factors are mentioned, we assume the potential outcome \$Y\$ depends only on \$S\$ and \$T\$. The distribution of stone sizes \$S\$ is independent of the treatment assignment mechanism in this hypothetical setup (or at least, the counterfactual distribution of \$S\$ for the treated group remains the same as the observed distribution of \$S\$ in the treated group, which is the standard assumption for calculating the ATE on the treated unless selection bias is explicitly defined by a confounder other than \$S\$).} \\
\\ 
\texttt{The counterfactual probability is:} \\
\texttt{\$\$P(Y=1 | T=1, do(T=0)) = P(Y=1 | S=\textbackslash{}text\{Small\}, T=} \\
            }
        }
    }
    & & \\ \\

    \theutterance \stepcounter{utterance}  
    & & & \multicolumn{2}{p{0.3\linewidth}}{
        \cellcolor[rgb]{0.9,0.9,0.9}{
            \makecell[{{p{\linewidth}}}]{
                \texttt{\tiny{[GM$|$GM]}}
                \texttt{No} \\
            }
        }
    }
    & & \\ \\

    \theutterance \stepcounter{utterance}  
    & & & \multicolumn{2}{p{0.3\linewidth}}{
        \cellcolor[rgb]{0.9,0.9,0.9}{
            \makecell[{{p{\linewidth}}}]{
                \texttt{\tiny{[GM$|$GM]}}
                \texttt{no} \\
            }
        }
    }
    & & \\ \\

    \theutterance \stepcounter{utterance}  
    & & & \multicolumn{2}{p{0.3\linewidth}}{
        \cellcolor[rgb]{0.9,0.9,0.9}{
            \makecell[{{p{\linewidth}}}]{
                \texttt{\tiny{[GM$|$GM]}}
                \texttt{game\_result = WIN} \\
            }
        }
    }
    & & \\ \\

\end{supertabular}
}

\end{document}
